\clearpage
\setcounter{chapter}{2}
\chaptermark{波动光学引论}
\pdfbookmark[0]{第二章 波动光学引论}{chapter2}
\problem{2-1}{
钠黄光系双线结构，其包含的两种谱线的波长分别为 $\lambda_1=5890\,\text{\AA}$、$\lambda_2=5896\,\text{\AA}$。设 $t=0$ 时刻，沿传播方向 $z$ 轴的原点 $O$ 处两波列的波峰重合。
\begin{enumerate}
\item 沿 $z$ 轴考察，两波列的波峰再次重合的距离 $z_0$ 为多远？
\item 若某光源发射双谱线 $\lambda'_1=570.0\,\mathrm{nm}$ 和 $\lambda'_2=600.0\,\mathrm{nm}$，那么上述要求下的纵向距离 $z'_0$ 为多少？
\item 从 $z_0$、$z'_0$ 的数值比较中可以引出怎样的定性结论？
\end{enumerate}
}
\problem{2-2}{
一束平行光其波长为 $\lambda$，其传播方向相对我们设定的坐标架 $(xyz)$ 的方向余弦角为 $(\alpha,\beta,\gamma)$，且 $\alpha=30^\circ$、$\beta=75^\circ$。
\begin{enumerate}
\item 试写出其波函数 $\widetilde U(x,y,z)$，设振幅为 $A$、原点相位 $\varphi_0=0$。
\item 试写出其波前函数 $\widetilde U(x,y)$。
\item 若方向角 $\beta$ 改变为 $\beta'=90^\circ$、$150^\circ$，试分别写出其波前函数 $\widetilde U_1(x,y)$ 和 $\widetilde U_2(x,y)$。
\end{enumerate}
}
\problem{2-3}{
如图所示，在一薄透镜的物方焦面上开有两个小孔 $O$ 和 $Q$ 而成为两个次波点源，且 $Q$ 点满足傍轴条件即 $a\ll F$（焦距）。
\begin{enumerate}
\item 试写出点源 $O$ 和 $Q$ 在透镜前表面 $(xy)$ 上产生的波前函数 $\widetilde U_1(x,y)$ 和 $\widetilde U_2(x,y)$。设振幅 $A_1=A_2=A$。
\item 经透镜以后，上述两个球面波成为两列平面波，而射向像方焦面 $(x'y')$，试写出各自的波前函数 $\widetilde U'_1(x',y')$ 和 $\widetilde U'_2(x',y')$。设振幅 $A'_1=A'_2=A$。
\end{enumerate}
\tikzfig{fig_2_3}{2-3}
}
\problem{2-4}{
约定光路自左向右。发射平面 $(x_0y_0)$ 在接收平面 $(xy)$ 左侧，两者纵向距离 $z=10^3\,\mathrm{mm}$。在发射平面上有三个点源，其横向位置坐标 $(x_0,y_0)$ 分别为 $(0,0)$、$(3\,\mathrm{mm},4\,\mathrm{mm})$ 和 $(-3\,\mathrm{mm},-4\,\mathrm{mm})$，且同相位。
\begin{enumerate}
\item 试分别写出这三个点源所产生的波前函数 $\widetilde U_1(x,y)$、$\widetilde U_2(x,y)$ 和 $\widetilde U_3(x,y)$。设其波数均为 $k$，振幅均为 $A$。要求作图示意。
\item 试写出三者各自的共轭波前 $\widetilde U_1^*(x,y)$、$\widetilde U_2^*(x,y)$ 和 $\widetilde U_3^*(x,y)$，并从中判断其所代表的光波的聚散性及中心坐标 $(x_0,y_0,z_0)$。
\end{enumerate}
}
\problem{2-5}{
在波前 $(xy)$ 平面上分别出现以下几种相位分布函数，
\begin{align*}
\widetilde U_1&\propto \exp\left(\mathrm{i}5k\frac{x^2+y^2}{D}\right),\\
\widetilde U_2&\propto \exp\left(\mathrm{i}k\frac{x^2+y^2}{2D}\right)\exp\left(-\mathrm{i}k\frac{5x+8y}{2D}\right),\\
\widetilde U_3&\propto \exp\left(-\mathrm{i}4k\frac{x^2+y^2}{2D}\right)\exp\left(-\mathrm{i}k\frac{5x+8y}{2D}\right).
\end{align*}
$k$ 为波数，$D>0$。试分别判断这些波前所联系的波场的类型和特征。
}
\problem{2-6}{
太阳看起来像个圆盘，这说明对地球上的观察者而言，太阳显然不是一个点源。那么，
\begin{enumerate}
\item 太阳上一点源发射的球面波，到达地球上在多大范围内可被看作一平面波。已知日地距离约为 $1.5\times10^8\,\mathrm{km}$，波长取 $0.5\,\mu\mathrm m$。
\item 对于月亮，其上一点源发射的球面波，在地球多大范围内可被看作一平面波。已知月地距离约为 $3.8\times10^5\,\mathrm{km}$。
\end{enumerate}
}
\problem{2-7}{
一射电源离地面高度 $h$ 约 $300\,\mathrm{km}$，向地面发射波长 $\lambda$ 约 $20\,\mathrm{cm}$ 的微波，接收器的孔径 $\rho$ 为 $2\,\mathrm m$。试问针对此孔径 $\rho$，这射电源的高度 $h$ 是否满足远场条件。
}
\problem{2-8}{
在杨氏双孔实验中，孔距为 $0.4\,\mathrm{mm}$，孔与接收屏的距离为 $3\,\mathrm m$，试分别对下列三条典型谱线求出干涉条纹的间距：
\[
\text{F 蓝线 }\lambda_1=4861\,\text{\AA},\quad
\text{D 黄线 }\lambda_2=5893\,\text{\AA},\quad
\text{A 红线 }\lambda_3=6563\,\text{\AA}.
\]
}
\problem{2-9}{
在杨氏双孔实验中，孔距为 $0.45\,\mathrm{mm}$，孔与接收屏的距离为 $1.2\,\mathrm m$，在某一准单色光照明下，测得 $10$ 条亮纹之间的距离为 $15\,\mathrm{mm}$。求光波长。
}
\problem{2-10}{
如图所示，两束相干平行光束其传播方向均平行于 $(xz)$ 面，对称地斜入射于记录介质 $(xy)$ 面上，光波长为氦氖激光的 $6328\,\text{\AA}$。
\begin{enumerate}
\item 当两束光之夹角为 $10^\circ$ 时，求干涉条纹的间距 $\Delta x$ 及相应的空间频率 $f$。
\item 当两束光之夹角为 $60^\circ$ 时，求干涉条纹的间距及相应的空间频率。
\item 若记录介质的空间分辨率为 $1500$ 线/$\mathrm{mm}$，试问这介质能否精确地记录下上述两种条纹。
\end{enumerate}
\tikzfig{fig_2_10}{2-10}
}
\problem{2-11}{
（接上题）若那两束平行光系来自同一光源所发射的自然光，且光强相等。
\begin{enumerate}
\item 当两束光之夹角 $\alpha=20^\circ$，求干涉场的衬比度 $\gamma$。
\item 当两束光之夹角 $\alpha=90^\circ$，求干涉场的衬比度 $\gamma$。
\item 若两束光强不相等，设 $I_2=mI_1$，试导出 $\gamma(\alpha,m)$ 函数形式。
\end{enumerate}
}
\problem{2-12}{
如图所示，三束完全相干的平行光投射于屏幕 $(xy)$，设其振幅为 $A_1$、$A_0=2A_1$、$A_2=A_1$；其初相位在原点均为 $0$。试分别采用复数法和矢量图解法，求出干涉场的复振幅分布 $\widetilde U(x,y)$，并据此讨论干涉场的主要特征，本题要求作图显示 $I(x)$ 曲线。
\tikzfig{fig_2_12}{2-12}
}
\problem{2-13}{
让我们来研究一平面波和一球面波的干涉场，这可通过一点源被置于一旋转抛物面镜焦点处来实现，参见题图。设平面波的振幅为 $A_1$，傍轴球面波到达记录介质平面 $(xy)$ 的振幅为 $A_2$，其发散中心到记录介质平面之距离为 $a$。
\begin{enumerate}
\item 试导出干涉场的波前函数 $\widetilde U(x,y)$ 和光强分布 $I(x,y)$。
\item 明示干涉花样的特征。
\end{enumerate}
\tikzfig{fig_2_13}{2-13}
}
\problem{2-14}{
讨论一个随机干涉场问题。

设想杨氏实验中的双孔即两个相干点源的间距 $d$ 在作无规则的颤动，而成为一个随机量 $\widetilde d=d_0+\widetilde\Delta$，这里 $d_0$ 是平均间距，$\widetilde\Delta$ 是一个随机涨落的量，且 $d_0\gg|\widetilde\Delta|\gg\lambda$；还有，这无规则颤动的周期远小于探测器的响应时间；当然，也无妨再设定 $\Delta$ 取值的几率分布函数 $f(\Delta)$ 服从某一统计分布律，比如高斯分布 $f(\Delta)\propto \mathrm e^{-a\Delta^2}$。试定性地描绘出所观测到的平均干涉强度曲线 $I(x)$。

提示：相位差
\[
\delta(x)=k\frac{d}{D}x=k\frac{d_0}{D}x+k\frac{\widetilde\Delta}{D}x;
\]
$\delta(x)$ 的无规性是受控于瞬时场点位置 $x$ 所控制。
\tikzfig{fig_2_14}{2-14}
}
\problem{2-15}{
在菲涅耳圆孔衍射实验中，点光源距离圆孔 $1.5\,\mathrm m$，接收屏距离圆孔 $6.0\,\mathrm m$，圆孔半径 $\rho$ 从 $0.50\,\mathrm{mm}$ 开始逐渐扩大，设光波长为 $0.63\,\mu\mathrm m$。求
\begin{enumerate}
\item 最先两次出现中心亮斑时圆孔的半径 $\rho_1$、$\rho_2$。
\item 最先两次出现中心暗斑时圆孔的半径 $\rho'_1$、$\rho'_2$。
\end{enumerate}
}
\problem{2-16}{
在菲涅耳圆孔衍射实验中，点光源距离圆孔 $2.0\,\mathrm m$，圆孔半径固定为 $2.0\,\mathrm{mm}$，波长为 $0.5\,\mu\mathrm m$。当接收屏由很远处向圆孔靠近时，
\begin{enumerate}
\item 前三次出现中心亮斑时的屏幕位置（与圆孔的距离）。
\item 前三次出现中心暗斑时的屏幕位置。
\end{enumerate}
}
\problem{2-17}{
用一直边刀片将点光源产生的波前遮住一半，问在一定距离的屏幕上几何阴影边缘点的衍射光强为多少（与自由传播光强 $I_0$ 比较）。
}
\problem{2-18}{
如图所示系 $6$ 个不同样式的衍射屏用以面对一平面光波，图旁的符号表示该波处到轴上观察点 $P$ 的距离，而 $b$ 正是这衍射屏中心到 $P$ 之距离。试分别给出衍射光强 $I(P)$，与自由传播光强 $I_0$ 相比较。
\tikzfig{fig_2_18}{2-18}
}
\problem{2-19}{
一菲涅耳波带片其第一个半波带的半径 $\rho_1$ 为 $5.0\,\mathrm{mm}$。
\begin{enumerate}
\item 若用波长为 $1.06\,\mu\mathrm m$ 的单色平行光照明，求其主焦距 $f$。
\item 若要求对此波长其主焦距缩短为 $25\,\mathrm{cm}$，需将此波带片精缩多少？
\end{enumerate}
}
\problem{2-20}{
一菲涅耳波带片对 $900\,\mathrm{nm}$ 的红外光其主焦距为 $30\,\mathrm{cm}$，若改用 $633\,\mathrm{nm}$ 的氦氖激光照明，其主焦距变为多少？
}
\problem{2-21}{
现手边有一张浮雕型全透明的菲涅耳波带片，粗测其对台灯灯光的主焦距为 $10\,\mathrm{cm}$，其有效半径为 $4\,\mathrm{cm}$。试估算该波带片聚光倍率的数量级（与自由光强相比较），可设波长为 $550\,\mathrm{nm}$，并忽略倾斜因子的影响。
}
\problem{2-22}{
试导出：
\begin{enumerate}
\item 平行光斜入射时单缝夫琅禾衍射强度公式为
\[
I(\theta)=I_0\left(\frac{\sin\alpha'}{\alpha'}\right)^2,\qquad
\alpha'=\frac{\pi a}{\lambda}(\sin\theta-\sin\theta_0),
\]
这里，$a$ 为单缝宽度，$\theta_0$ 为入射光束的倾角，$\theta$ 为衍射角，两者均相对于单缝屏的法线方向。
\item 在此斜入射条件下，零级衍射斑的半角宽度公式为
\[
\Delta\theta_0=\frac{\lambda}{a\cos\theta_0}.
\]
\end{enumerate}
}
\problem{2-23}{
考虑到介质界面宽度为有限值，反射光束和折射光束必然有一定的衍射发散角 $(\Delta\theta,\Delta\theta')$。试估算在以下不同入射角 $i$ 时的 $(\Delta\theta,\Delta\theta')$ 值。设界面宽度 $a$ 为 $1\,\mathrm{cm}$，光波长 $\lambda_0$ 为 $600\,\mathrm{nm}$，介质折射率 $n$ 为 $1.5$。
\begin{enumerate}
\item $i\approx0^\circ$ 即正入射；
\item $i\approx75^\circ$；
\item $i\approx88^\circ$ 即掠入射。
\end{enumerate}
}
\problem{2-24}{
氦氖激光器的发光区集中于一毛细管，其管径约 $2\,\mathrm{mm}$。
\begin{enumerate}
\item 试估算从管口端面出射的 He-Ne 激光束，其衍射发散角为多大，设波长为 $633\,\mathrm{nm}$。
\item 若此光束射至 $10\,\mathrm m$ 远的屏幕上，其光斑尺寸为多大？
\item 若此光束射至月球表面，其光斑尺寸为多大？
\end{enumerate}
}
\problem{2-25}{
考虑衍射色散效应。在单缝夫琅禾衍射实验中，入射光含有两种波长，蓝光 $\lambda_1\approx400\,\mathrm{nm}$ 和红光 $\lambda_2\approx700\,\mathrm{nm}$，两者的零级斑中心自然是重合的，设两者的入射强度相等。试问：
\begin{enumerate}
\item 两者在零级斑中心点的光强是否相等？如否，其比值为多少？
\item 两者零级斑尺寸比值为多少？
\end{enumerate}
}
\problem{2-26}{
一衍射细丝测径仪，它将单缝夫琅禾衍射装置中的单缝替换为细丝，今测得其产生的零级衍射斑的宽度即两个一级暗点间的距离为 $1\,\mathrm{cm}$，求细丝的直径 $a$。已知光波长 $633\,\mathrm{nm}$，透镜焦距为 $50\,\mathrm{cm}$。
}
\problem{2-27}{
在单缝或单孔夫琅禾衍射实验中，若其装置有如下几种变动时，试讨论相应的衍射图样将有怎样的变化：
\begin{enumerate}
\item 增大接收透镜的焦距，
\item 增大接收透镜的口径，
\item 衍射屏沿系统的纵向作前后平移，
\item 衍射屏沿横向平移，
\item 衍射屏绕纵向轴而旋转。
\end{enumerate}
}
\problem{2-28}{
一对双星的角间隔为 $0.05''$，试问：
\begin{enumerate}
\item 需要至少多大口径的望远镜才能分辨它俩？
\item 与此口径相匹配的望远镜的角放大率应当设计为多少？
\end{enumerate}
}
\problem{2-29}{
一台显微镜其数值孔径 $\mathrm{N.A.}\approx1.32$，物镜焦距 $f_o$ 为 $1.91\,\mathrm{mm}$，而目镜焦距 $f_e$ 为 $50\,\mathrm{mm}$。试求：
\begin{enumerate}
\item 其最小分辨间隔 $\delta y_m$。
\item 其有效放大率 $M_{\mathrm{eff}}$。
\item 该显微镜的光学筒长约为多少？
\end{enumerate}
}
\problem{2-30}{
一照相机在离地面 $200\,\mathrm{km}$ 的高空拍摄地面上的物体，若要求它能分辨地面上相距 $1\,\mathrm m$ 的两点，问：
\begin{enumerate}
\item 此照相机的镜头至少需要多大？设镜头的几何像差已很好地消除，感光波长为 $400\,\mathrm{nm}$。
\item 与之匹配的感光胶片其分辨率至少应当为多少？设该镜头的焦距为 $8\,\mathrm{cm}$。
\end{enumerate}
}
\problem{2-31}{
用口径为 $1\,\mathrm m$ 的光学望远镜，能分辨月球表面上两点的最小距离为多少？已知地月距离约为 $3.8\times10^5\,\mathrm{km}$。
}
\problem{2-32}{
在水下有一超声探测器，其圆形孔径为 $60\,\mathrm{cm}$，发射 $40\,\mathrm{kHz}$ 的超声波，设声速为 $1.5\,\mathrm{km/s}$。问：
\begin{enumerate}
\item 此超声束的发散角 $\Delta\theta$ 为多少？
\item 在距离 $1\,\mathrm{km}$ 远该超声波照射的范围有多大？
\end{enumerate}
}
\problem{2-33}{
一束自然光正入射于重叠的两张偏振片上，如果透射光强为
\begin{enumerate}
\item 入射光强的 $1/3$；
\item 透射光束最大光强的 $1/3$，
\end{enumerate}
试分别确定那两个偏振片的透振方向之夹角 $\theta$。
}
\problem{2-34}{
在一对正交偏振片之间插入另一张偏振片，其透振方向沿 $45^\circ$ 角（相对那一对正交的透振方向），当自然光入射时，求透射光强的百分比（相对于入射光强）。
}
\problem{2-35}{
一束自然光正入射于一组合有 $4$ 张的偏振片，其每片的透振方向相对于前面一片均沿顺时针方向转过 $30^\circ$ 角，求最终透射光强的百分比（相对于入射光强）。
}
\problem{2-36}{
要使一束线偏振光的振动方向转过 $90^\circ$，且要求最终透射光强为原来入射光强的 $95\%$，试问至少需要多少块理想偏振片？

提示：透振方向依次转过相同角度 $\alpha$。
}
\problem{2-37}{
一张偏振片正对着一束部分偏振光，当它相对光强极大值 $I_M$ 的方位转过 $45^\circ$ 时，透射光强为 $I_M$ 的 $2/3$。求入射光的偏振度 $p$。
}
\problem{2-38}{
一偏振片正对着一光束，其包含两个非相干的线偏振光，两者光强相等均为 $I_0$，偏振方向之夹角 $\alpha_0$ 为 $60^\circ$，参见题图。 当偏振片旋转一周时，
\begin{enumerate}
\item 出现几次光强极大和光强极小。
\item 求出光强极大值及其方位角 $(I_M,\alpha_M)$，光强极小值及其方位角 $(I_m,\alpha_m)$。以 $x$ 轴为参考来标定 $\alpha$ 角。
\item 若两者光强不相等，比如 $I_2=2I_1$，其他条件不变，试求出 $(I_M,\alpha_M)$ 和 $(I_m,\alpha_m)$。
\end{enumerate}
\tikzfig{fig_2_38}{2-38}
}
\problem{2-39}{
光场中某一点的光矢量 $\boldsymbol E(t)$ 在其横平面上的两个正交分量为（这里不计较单位），
\[
E_x(t)=5\cos\omega t,\qquad E_y(t)=5\cos(\omega t+60^\circ).
\]
现用一偏振片正对着这束光且旋转之，以考察其透射光强的变化。你认为其合成光矢量是一圆偏振光吗？如非，试给出其光强极大值、极小值及相应的方位角，即给出 $(I_M,\alpha_M)$、$(I_m,\alpha_m)$。
}
\problem{2-40}{
光场中某一点的光矢量 $\boldsymbol E(t)$ 在其横平面上的两个正交分量为（这里不计较单位），
\[
E_x(t)=3\cos\omega t,\qquad E_y(t)=4\cos(\omega t+\delta).
\]
现用一偏振片面对这束光且旋转之，以考察其透射光强的变化。试就下面三种情况分别回答光强极大值及其方位角 $(I_M,\alpha_M)$，光强极小值及其方位角 $(I_m,\alpha_m)$，以 $x$ 轴为参考来标定 $\alpha$ 角。
\begin{enumerate}
\item $\delta=0$，
\item $\delta=\pm\pi/2$，
\item $\delta=\pi/6$ 即 $30^\circ$。
\end{enumerate}
}
\problem{2-41}{
一张偏振片正对着一光束，其包含两个非相干的椭圆偏振光，两者光强相等均为 $I_0$，且长短轴之比也相同，只是长短轴取向互换，参见题图。试证明，在这种情形下偏振片转动过程中，透射光强保持为一常量 $I_0$。

注：这一结论将在第 8 章 8.4 节为了区分自然光与圆偏振光时用到。
\tikzfig{fig_2_41}{2-41}
}
\problem{2-42}{
光强为 $I_0$ 的自然光相继通过三个偏振片 $P_1$、$P_2$ 和 $P_3$，其中 $P_1$ 与 $P_3$ 静止且其透振方向彼此正交，而 $P_2$ 以角速度 $\omega_0$ 绕光线为轴旋转。试求最终透射光强 $I_3(t)$。可设 $t=0$ 时刻，$P_2$ 透振方向平行 $P_1$；光扰动频率为 $\omega$。
}
